\chapter{はじめに}

このガイドは、\textbf{マニュアル第1弾（WSL 入門）の続編}です。
WSL（Ubuntu）のインストールとターミナルの基本操作は完了している前提で進めます。

\section{このガイドでやること}

\begin{itemize}
  \item 仮想環境とは何か、なぜ必要かを理解する
  \item パッケージマネージャ \textbf{uv} をインストールする
  \item \cmd{git clone} で課題リポジトリを取得する
  \item \cmd{uv sync} で必要なライブラリをまとめてインストールする
  \item VS Code でノートブックを開いて課題を実行する
\end{itemize}

\section{ゴール}

このガイドを読み終えると、VS Code で課題フォルダを開くだけで
作業を始められるようになります。

\begin{wslbox}[毎回の起動（最終目標）]
\begin{lstlisting}[style=bash]
$ code ~/NMR-lab-kadai/2026年度_B4課題引継ぎ
\end{lstlisting}
\end{wslbox}

\section{前提条件}

\begin{itemize}
  \item WSL（Ubuntu）がインストール済みであること
  \item WSL のターミナルを開けること
  \item インターネットに接続できること（初回セットアップ時のみ）
  \item VS Code に Jupyter 拡張機能がインストールされていること
\end{itemize}

\begin{notebox}
\textbf{Python のインストールは不要です。}\\
uv が Python 本体も含めて自動でダウンロード・管理します。
\end{notebox}
